\documentclass[aspectratio=169]{beamer}
\usetheme{Madrid}
\usecolortheme{whale}

\usepackage[utf8]{inputenc}
\usepackage[spanish]{babel}
\usepackage{amsmath,amssymb}
\usepackage{graphicx}
\usepackage{listings}
\usepackage{booktabs}

\lstset{
    language=Python,
    basicstyle=\ttfamily\footnotesize,
    keywordstyle=\color{blue}\bfseries,
    stringstyle=\color{red},
    commentstyle=\color{gray}\itshape,
    breaklines=true,
    showstringspaces=false
}

\title[INT210 - Sesión 0.6]{Sesión 0.6: Diccionarios, \texttt{defaultdict} y Frecuencias}
\subtitle{Módulo 0.2: Estructuras de Datos Nativas en Python}
\author[Diplomado en Ciencia de Datos]{Cátedra de Inteligencia Artificial y Ciencia de Datos}
\institute[INT210]{Machine Learning From Scratch}
\date{Semana 2}

\begin{document}

\frame{\titlepage}

\begin{frame}{Contenido de la Sesión}
    \tableofcontents
\end{frame}

\section{El Cofre de Llaves y Tablas Hash}

\begin{frame}{La Analogía: El Cofre de Llaves Mágicas y la Libreta}
    \begin{columns}
        \column{0.55\textwidth}
        \begin{itemize}
            \item \textbf{Diccionario (\texttt{dict}):} Llaves mágicas que abren cajas en $\mathcal{O}(1)$. Si la llave no existe, genera \texttt{KeyError}.
            \item \textbf{\texttt{defaultdict(int)}:} Libreta con auto-reposición. Si una clave no existe, la inicializa en $0$ silenciosamente.
            \item \textbf{\texttt{Counter}:} Especialista en conteo de frecuencias y extracción del top $k$ de mayor densidad.
        \end{itemize}
        \column{0.45\textwidth}
        \begin{alertblock}{Principio de Complejidad}
            El acceso hash directo $\mathcal{O}(1)$ es la base del Procesamiento de Lenguaje Natural (NLP) y de la indexación de eventos de seguridad.
        \end{alertblock}
    \end{columns}
\end{frame}

\section{Estructura Interna y Mitigación de KeyError}

\begin{frame}[fragile]{Evolución de Técnicas de Conteo}
    \textbf{1. Conteo Seguro con \texttt{.get()}:}
\begin{lstlisting}
conteo[palabra] = conteo.get(palabra, 0) + 1
\end{lstlisting}
    \vspace{0.2cm}
    \textbf{2. Conteo con \texttt{defaultdict(int)}:}
\begin{lstlisting}
conteo_default = defaultdict(int)
for palabra in tokens:
    conteo_default[palabra] += 1
\end{lstlisting}
    \vspace{0.2cm}
    \textbf{3. Extracción de los $k$ Más Frecuentes con \texttt{Counter}:}
\begin{lstlisting}
top_amenazas = Counter(tokens).most_common(5)
\end{lstlisting}
\end{frame}

\section{Aplicación en Ciberseguridad y NLP}

\begin{frame}{Bolsa de Palabras (*Bag-of-Words*) para Detección de Phishing}
    Cada mensaje se proyecta como una combinación lineal de frecuencias léxicas:
    \begin{equation*}
        \text{Score}(\text{mensaje}) = \sum_{w \in \text{mensaje}} \text{freq}(w) \cdot \text{PesoRiesgo}(w)
    \end{equation*}
    \begin{block}{Regla de Decisión}
        \begin{equation*}
            \text{Clase} = \begin{cases}
                \text{Phishing}, & \text{si } \text{Score} \ge \theta \\
                \text{Legítimo}, & \text{si } \text{Score} < \theta
            \end{cases}
        \end{equation*}
    \end{block}
\end{frame}

\begin{frame}{Conclusiones}
    \begin{itemize}
        \item Las tablas hash nativas son la estructura óptima para indexación de texto y métricas discretas.
        \item \texttt{defaultdict} y \texttt{Counter} simplifican drásticamente el código y previenen excepciones en producción.
        \item \textbf{Siguiente Sesión (0.7):} Sets y Álgebra de Conjuntos en Alta Velocidad.
    \end{itemize}
\end{frame}

\end{document}
